\section{Discussion}
\label{sec:discussion}

\subsection{The Peat-Drainage Nexus and Hydrological Vulnerability}
\label{subsec:disc_nexus}

Our spatial and TreeSHAP findings provide direct empirical confirmation of the \textit{peat-drainage nexus} operating across degraded tropical peatlands \citep{Hooijer2012Biogeosciences, Ritzema2014HESS}. In pristine equatorial peat swamp forests, high water tables (typically within 10--20\,cm of the surface) and dense forest canopies maintain water saturation throughout the acrotelm, rendering peat domes virtually immune to deep fire ignition even during seasonal dry periods \citep{Wosten2008Geoderma, Dadap2021NatComm}. 

The excavation of extensive primary and secondary drainage canals across the Mega Rice Project (Blocks A, B, and C) permanently severed this self-regulating hydrological mechanism \citep{Silvius2005Wetlands, Ritzema2014HESS}. By piercing the elevated crests of ombrotrophic peat domes, canals induce continuous gravitational drawdown of phreatic water tables toward regional river systems. Our data reveal that \textbf{74.3\%} of all active fire hotspots in Central Kalimantan ignited within delineated peatlands, with over \textbf{52.5\%} situated in deep ($>2$\,m) and very deep ($>3$\,m) peat domes. Within the ex-PLG landscape, metric Euclidean distance to primary canals emerged as the second most influential predictor of severe Fire Radiative Power ($\text{mean } |\text{SHAP}| = 0.0322$).

Canal drainage induces severe physical compaction, irreversible desiccation, and oxidative loss of the fibrous peat matrix \citep{Hooijer2012Biogeosciences, Graham2022SciRep}. When the water table drops below $-40$\,cm, aerobic microbial decomposition accelerates, transforming the uppermost meter of peat into a dry, porous fuel bed with exceptionally low bulk density and high air-filled porosity \citep{Usup2004Tropics}. Once ignited, combustion shifts from brief above-ground surface flaming to prolonged, oxygen-limited subterranean smoldering throughout the catotelm \citep{Wooster2018NatGeosci, Vetrita2020RemoteSensing}. The intense thermal radiation released along canal corridors confirms that artificial drainage vectors act as regional fire vulnerability epicenters, requiring prioritized hydrological restoration.

\subsection{Resolving the Land Tenure Governance Void in Degraded Peatlands}
\label{subsec:disc_tenure}

Our cadastral spatial overlay resolves a long-standing point of contention in environmental policy by quantifying the land tenure attribution of tropical peat fires: \textbf{75.42\% of all verified active fire hotspots (33,138 detections) occurred on unallocated Non-Concession land (\textit{Area Penggunaan Lain} / APL)}, while certified oil palm concessions (HGU) and commercial forestry licenses (PBPH) accounted for \textbf{16.56\%} and \textbf{8.02\%}, respectively.

These findings decisively challenge simplistic single-sector narratives that attribute Indonesian peat wildfires almost exclusively to corporate concession clearing \citep{Tacconi2003CIFOR, Carlson2012PNAS}. This empirical reality does not absolve corporate concessionaires: \textbf{24.58\% of all active fires (10,799 hotspots)} ignited directly within legal corporate boundaries. Under Indonesian environmental jurisprudence (Law No.~32/2009 on Environmental Protection and Management, Article 88), concession holders bear strict liability (\textit{tanggung jawab mutlak}) for any fires occurring within their demarcated concessions, regardless of fault. The spatial identification of specific entities experiencing repeated fire clusters (e.g., PT.~Sangkowong Sinta with 954 hotspots; PT.~Rimba Makmur Utama with 964 hotspots) provides transparent, auditable evidence for corporate Environmental, Social, and Governance (ESG) compliance monitoring and legal enforcement.

However, the overwhelming concentration of fires in unallocated APL zones exposes an institutional and administrative vacuum across Central Kalimantan's degraded peatlands \citep{Cattau2016GCB, Putranto2023LandUsePol}. Following the formal abandonment of the Mega Rice Project in 1999, hundreds of thousands of hectares of drained peatland were designated as state land without assigned institutional managers or active concessionaires. In the absence of state enforcement, these lands have evolved into de facto open-access regimes characterized by speculative land acquisition, opportunistic clearance along canal banks, and accidental ignitions by informal resource extractors \citep{Cattau2016GCB}. Because no corporate entity is legally liable for fire suppression across unallocated APL lands, fire-fighting responsibilities fall on under-resourced regional disaster management agencies (BPBD), forest fire brigades (Manggala Agni), and volunteer community groups (\textit{Masyarakat Peduli Api} / MPA). Without an institutional governance framework to enclose and manage these unallocated APL peatlands, regulatory pressure on corporate concessions alone will remain fundamentally incapable of halting regional wildfire disasters.

\subsection{Methodological Imperatives in Spatial Wildfire Machine Learning}
\label{subsec:disc_spatial_ml}

The comparative evaluation between Model~A and Model~B provides a critical methodological lesson for geospatial data science, remote sensing, and disaster modeling. In published environmental machine learning literature, the inclusion of geodetic coordinates (\texttt{latitude} and \texttt{longitude}) as independent predictors is widely practiced under the presumption that spatial coordinates capture unobserved geographical heterogeneity \citep{Meyer2018EcolModel, Dormann2007Ecography}.

Our results demonstrate that this practice introduces a dangerous illusion of predictive skill. In Model~A, spatial coordinates acted as algorithmic shortcuts: decision trees recursively split coordinate axes to isolate historical spatial clusters. When evaluated using standard random train-test splitting (75:25), Model~A achieved an ostensibly superior ROC-AUC of 0.7956 because testing samples occupied the same spatial autocorrelation neighborhoods as training samples \citep{Tobler1970EconGeogr, Legendre1993Ecology}. However, when subjected to \textbf{Spatial Block Cross-Validation} (\texttt{GroupKFold} by regency), Model~A's predictive capacity collapsed: mean ROC-AUC plunged to 0.5399, and the $F_1$-score plummeted to 0.0583, collapsing to \textbf{0.0000 across four of five spatial folds}. The model was incapable of predicting severe fires in novel geographical regions because its learned decision boundaries relied on coordinate partitions that did not exist in out-of-sample blocks \citep{Roberts2017Ecography, Wadoux2021EcolModel}.

Conversely, by completely excluding geographic coordinates, Model~B forced the Random Forest ensemble to learn functional, physically grounded dependencies on antecedent precipitation deficits ($P_{\text{7d}}$), Euclidean canal proximity ($d_{\text{canal}}$), wind speed, and peat depth. While its random holdout ROC-AUC was slightly lower (0.7527), Model~B maintained positive $F_1$-scores across \textbf{all five spatial blocks}, achieving a mean spatial $F_1$-score of 0.1399---representing a **2.40-fold improvement** over Model~A. These findings demonstrate that machine learning models in environmental remote sensing must be validated via spatial block cross-validation and constrained to their valid Area of Applicability \citep{MeyerPebesma2021NatComm, Schratz2019EcolModel}.

\subsection{Study Limitations, Uncertainties, and Threats to Validity}
\label{subsec:disc_limitations}

To ensure complete transparency and uphold the highest standards of scientific rigor, several methodological limitations, data uncertainties, and threats to validity must be articulated:

\subsubsection{Sensor Resolution, Orbital Sampling Gaps, and Smoldering Combustion Physics}
First, active fire observations from the NASA FIRMS constellation (VIIRS 375\,m and MODIS 1\,km) are subject to inherent satellite remote sensing limitations \citep{Schroeder2021RSE, Giglio2020RSE}. Both sensor systems operate on sun-synchronous, polar-orbiting platforms that pass over Central Kalimantan only two to four times per day, with daytime overpasses concentrated between 12:30 and 14:30 local solar time. Consequently, diurnal fire front variations occurring outside orbital passes, as well as nocturnal fire spread, escape observation. Furthermore, persistent cloud cover and dense convective smoke plumes frequently attenuate or obstruct infrared radiance, introducing potential omission errors during peak burning periods \citep{Giglio2020RSE}.

More fundamentally, infrared active fire algorithms detect photon radiance emitted by high-temperature surface flaming combustion (typically $>800$\,K). In degraded tropical peatlands, however, destructive burning frequently occurs as subterranean smoldering combustion within the catotelm layer at temperatures of 500--700\,K without open flames \citep{Usup2004Tropics, Wooster2018NatGeosci, Vetrita2020RemoteSensing}. When smoldering burns beneath an intact forest canopy or deep within the peat matrix, emitted thermal energy is partially trapped or absorbed, reducing the surface Fire Radiative Power. Therefore, our target threshold ($\text{FRP} \ge 12.0$\,MW) models ``intense radiative thermal emission'' rather than the volumetric mass or subterranean depth of consumed peat.

\subsubsection{Scale Mismatch in Atmospheric Reanalysis and Potential Ecological Fallacy}
Second, meteorological forcing derived from ECMWF ERA5-Land reanalysis features a nominal spatial resolution of $0.1^\circ \approx 9$\,km, aggregated in our processing pipeline at the administrative regency level on a daily cadence. In extensive jurisdictions such as Kapuas Regency (spanning over 15,000\,km$^2$ from the northern highlands to the southern Java Sea coast), assigning uniform daily meteorological values introduces an inevitable scale mismatch and potential ecological fallacy. Localized convective thunderstorms that deposit heavy precipitation over an isolated peat dome cannot be resolved at this spatial scale. Future operational implementations should integrate dynamically downscaled weather radar data or dense automated weather station (AWS) networks.

\subsubsection{Temporal Scope of the Acute Fire Crisis}
Third, the machine learning models were trained and cross-validated on an acute eight-day fire crisis window (5--12 October 2026, comprising 43,937 observations) characterizing the transition from severe drought to monsoonal rainfall. While this sample captures extreme fire behavior under critical fuel desiccation conditions, the empirical parameter weights reflect this specific climatological pulse. Model application across the early dry season (June--July) or during prolonged multi-year wet regimes (e.g., strong La Ni\~na cycles) warrants dynamic seasonal retraining and domain adaptation.

\subsubsection{Spatial Transferability, Covariate Shift, and Class Imbalance Trade-offs}
Fourth, while Model~B completely eliminated catastrophic spatial memorization collapse, its mean Spatial Block CV ROC-AUC averaged 0.4950 ($\pm$0.0490) and its spatial CV $F_1$-score reached 0.1399. Rather than an algorithmic flaw, this outcome reflects substantial \textit{spatial covariate shift} between disparate administrative blocks. For example, primary canal drainage networks are concentrated in Pulang Pisau and Kapuas (ex-PLG Blocks A and B), whereas Seruyan Regency features distinct, uncanalized coastal peat domes with different hydrological boundary conditions. Furthermore, the extreme class imbalance (15.13\% severe fire class; 1:5.6 ratio) combined with inverse class weighting (\texttt{class\_weight=balanced}, $w_1 = 3.30$) intentionally prioritized recall over precision. On the independent holdout test set ($n = 10,985$), Model~B captured 65.82\% of all severe fires ($\text{TP} = 1,094$), but generated 2,737 false positives ($\text{FP}$), yielding an extreme-class precision of 28.56\% and an overall holdout $F_1$-score of 0.3983.

In wildfire disaster management, this trade-off is operationally justified by \textbf{loss asymmetry}: a false negative (failing to detect a severe fire, allowing it to penetrate the subterranean peat column) incurs catastrophic ecological, health, and economic costs. Conversely, a false positive triggers precautionary aerial or ground reconnaissance. Nevertheless, excessive false alarms risk inducing \textbf{alert fatigue} among emergency response personnel. To operationalize the model safely, we recommend implementing a \textbf{Multi-Tier Risk Stratification Protocol} based on continuous predicted probabilities:
\begin{itemize}
    \item \textit{Advisory / Waspada} ($p \in [0.30, 0.50)$): Heightened community patrolling and water table monitoring.
    \item \textit{Watch / Siaga} ($p \in [0.50, 0.70)$): Ground dispatch of Manggala Agni teams to test surface peat moisture.
    \item \textit{Warning / Awas} ($p \ge 0.70$): Immediate deployment of aerial water-bombing helicopters and activation of canal-blocking control gates.
\end{itemize}

\subsubsection{Cadastral Boundary Resolution and APL Land Tenure Heterogeneity}
Finally, the attribution of 75.42\% of active fires to unallocated Non-Concession (APL) land represents an aggregate legal classification. In practice, APL land in Central Kalimantan encompasses a diverse mosaic: customary (\textit{adat}) ancestral claims, transmigration farming settlements, independent smallholder oil palm plots with formal land titles (\textit{Sertifikat Hak Milik} / SHM), and abandoned, degraded state peatlands \citep{Putranto2023LandUsePol}. Disentangling these heterogeneous sub-categories requires fine-scale participatory cadastral mapping and village-level tenure boundary demarcation.
